\definecolor{promptbg}{gray}{0.95}
\definecolor{promptframe}{gray}{0.5}
\definecolor{afterbg}{HTML}{F0FAF0}
\definecolor{afterframe}{HTML}{2E7D32}
\definecolor{beforebg}{HTML}{FFF5F5}
\definecolor{beforeframe}{HTML}{C62828}
\definecolor{thinkbg}{HTML}{FAFAFA}
\definecolor{thinkframe}{gray}{0.75}
% --- mdframed style definitions (add to preamble) ---
\mdfdefinestyle{promptbox}{%
  backgroundcolor=promptbg,
  linecolor=promptframe,
  linewidth=0.8pt,
  innertopmargin=8pt,
  innerbottommargin=8pt,
  innerleftmargin=8pt,
  innerrightmargin=8pt,
  skipabove=6pt,
  skipbelow=4pt
}
\mdfdefinestyle{afterbox}{%
  backgroundcolor=afterbg,
  linecolor=afterframe,
  linewidth=1pt,
  innertopmargin=8pt,
  innerbottommargin=8pt,
  innerleftmargin=8pt,
  innerrightmargin=8pt,
  skipabove=6pt,
  skipbelow=4pt
}
\mdfdefinestyle{beforebox}{%
  backgroundcolor=beforebg,
  linecolor=beforeframe,
  linewidth=1pt,
  innertopmargin=8pt,
  innerbottommargin=8pt,
  innerleftmargin=8pt,
  innerrightmargin=8pt,
  skipabove=6pt,
  skipbelow=4pt
}
\mdfdefinestyle{thinkbox}{%
  backgroundcolor=thinkbg,
  linecolor=thinkframe,
  linewidth=0.5pt,
  innertopmargin=6pt,
  innerbottommargin=6pt,
  innerleftmargin=6pt,
  innerrightmargin=6pt,
  skipabove=4pt,
  skipbelow=4pt
}
% ============================================================
% Place the following in the body of the document
% ============================================================


\begin{figure*}[t]
\scriptsize
% === Full-width PROMPT ===
\begin{mdframed}[style=promptbox]
\noindent\textbf{프롬프트}
탄저병이라는 곰팡이성 질병에 대한 저항성에 기여하는 유전자를 찾기 위해 white lupine에서 high-throughput 실험을 수행한다. 기능이 알려지지 않은 세 개의 후보 유전자~-- G1, G2, G3를 받는다. 세 개의 knock-out 돌연변이체 (g1, g2, g3)와 이중 돌연변이체 (g1g2, g1g3, g2g3)를 만든다. 이 유전자들 가운데 적어도 하나는 다른 유전자(들)의 upstream에서 작동하는 transcription factor라는 것을 알고 있다. 병원체로 시험한 뒤 다음 결과를 받는다 (100\% = 야생형 저항성):

\vspace{2pt}
\begin{tabular}{@{}lr@{\hspace{6pt}}lr@{\hspace{6pt}}lr@{}}
\toprule
유전자형 & 저항성 & 유전자형 & 저항성 & 유전자형 & 저항성 \\
\midrule
g1 & 75\% & g1g3 & 10\% & g1g2 & 0\% \\
g2 & 0\%  & g2g3 & 0\%  &      &     \\
g3 & 50\% &      &      &      &     \\
\bottomrule
\end{tabular}
\vspace{2pt}

\noindent 어떤 결론을 도출할 수 있는가?

(A)~G2는 TF이고, G1\&G3는 pleiotropy를 보이며, G1은 G3에 대해 epistatic이다

(B)~G2는 TF이고, G1\&G3는 gene redundancy를 보이며, G1은 G3에 대해 epistatic이다

(C)~G2는 TF이고, G1\&G3는 promoter를 공유하며, G3는 G1에 대해 epistatic이다

(D)~G1은 TF이고, G2\&G3는 pleiotropy를 보이며, G2는 G1에 대해 epistatic이다
\hfill\textbf{\textit{정답:}~(B)}
\end{mdframed}
\vspace{2pt}
% === Two columns: Before \sdrlfootnotesize | After \sdrlfootnotesize ===
\noindent
\begin{minipage}[t]{0.49\textwidth}
% --- Before \sdrlfootnotesize ---
\begin{mdframed}[style=beforebox]
\noindent\textbf{\sdrlfootnotesize 이전} \hfill {\colorbox{beforebg}{\textcolor{beforeframe}{\textbf{$\boldsymbol\times$~오답 (비어 있음)}}}}

\noindent\textit{사고 과정 (일부 생략):}
\begin{mdframed}[style=thinkbox]
\begin{itemize}[leftmargin=12pt, labelsep=4pt, itemsep=0pt, parsep=0pt, topsep=0pt]
\item 야생형 (control): 100\% 저항성.
  \begin{itemize}[leftmargin=12pt, labelsep=4pt, itemsep=0pt, parsep=0pt, topsep=0pt]
  \item g1 돌연변이체: 75\% 저항성 (소폭 감소).
  \item g2 돌연변이체: 0\% 저항성 (완전 상실).
  \item g3 돌연변이체: 50\% 저항성 (중간 수준 감소).
  \item g1g3 이중 돌연변이체: 10\% 저항성 (유의미한 감소, g1 또는 g3 단독보다 더 낮음).
  \item g2g3 이중 돌연변이체: 0\% 저항성 (g2와 동일).
  \item g1g2 이중 돌연변이체: 0\% 저항성 (g2와 동일). [...]
  \item G2가 G1과 G3를 활성화하는 transcription factor (TF)라면:
    \begin{itemize}[leftmargin=12pt, labelsep=4pt, itemsep=0pt, parsep=0pt, topsep=0pt]
    \item G2가 G라면, G1과 G3는 표적이다.
    \item G2가 knock-out되면, G1과 G3는 발현되지 않는다.
    \item 따라서 g2는 g2g1 및 g2g3와 같은 표현형을 보여야 한다.
    \item 이는 데이터와 일치한다: g2 (0\%), g2g1 (0\%), g2g3 (0\%).
    \end{itemize}
  \item 이제 G1과 G3를 다시 보자:
    \begin{itemize}[leftmargin=12pt, labelsep=4pt, itemsep=0pt, parsep=0pt, topsep=0pt]
    \item g1: 75\%
    \item g3: 50\%
    \item g1g3: 10\%
    \item G1과 G3가 병렬 경로에서 작동한다면 ( pathways,, loss111111111111111111111 G11 G G11 G G G1 G G1 G G G G G G G1 G1 G G1 G G G G G G G1 G G G G G G111 G111 G G G1111 G111111111 G1111 G G1111111111111111111
    \end{itemize}
  \end{itemize}
\end{itemize}
{\itshape [\ldots 최대 token 예산에 도달할 때까지 숫자 1 token이 반복됨.]}
\end{mdframed}
\vspace{2pt}
\noindent\textit{응답:}
\begin{mdframed}[style=thinkbox]
\textit{[비어 있음 --- 응답이 생성되지 않음]}
\end{mdframed}
\end{mdframed}
\end{minipage}%
\hfill
\begin{minipage}[t]{0.49\textwidth}
% --- After \sdrlfootnotesize ---
\begin{mdframed}[style=afterbox]
\noindent\textbf{\sdrlfootnotesize 이후} \hfill {\colorbox{afterbg}{\textcolor{afterframe}{\textbf{\checkmark~정답}}}}

\noindent\textit{사고 과정 (일부 생략):}
\begin{mdframed}[style=thinkbox]
\begin{itemize}[leftmargin=12pt, labelsep=4pt, itemsep=0pt, parsep=0pt, topsep=0pt]
\item 야생형 (control): 100\% 저항성.
  \begin{itemize}[leftmargin=12pt, labelsep=4pt, itemsep=0pt, parsep=0pt, topsep=0pt]
  \item g1: 75\% 저항성 (약간 감소).
  \item g2: 0\% 저항성 (완전히 감수성).
  \item g3: 50\% 저항성 (중간 정도 감소).
  \item g1g3: 10\% 저항성 (매우 높은 감수성).
  \item g2g3: 0\% 저항성 (완전히 감수성).
  \item g1g2: 0\% 저항성 (완전히 감수성).
  \end{itemize}
\end{itemize}
g2는 0\% 저항성을 보인다. 이는 G2가 이 저항성에 필수적임을 뜻한다. g1g3는 10\%로, g1 (75\%)과 g3 (50\%)보다 훨씬 낮다. 이는 G1과 G3가 병렬 경로 (redundancy)에서 작동함을 시사한다.

\noindent G2가 G1과 G3의 upstream에 있는 TF라면, G2 (g2)를 knock-out하면 G1과 G3 모두의 발현이 막힌다. 이는 g2, g2g3, g1g2가 모두 0\% 저항성을 보이는 이유를 설명한다.

\vspace{2pt}
\noindent\textit{[...epistasis의 정의를 다시 검토하고, 모든 선택지를 확인한 뒤, (B)를 확정함...]}
\vspace{2pt}

\noindent 선택지 가운데 (B)만이 G2를 transcription factor로, G1/G3를 redundant한 관계로 올바르게 식별한다.
\end{mdframed}
\vspace{2pt}
\noindent\textit{응답:}
\begin{mdframed}[style=thinkbox]
유전자 G1, G2, G3에 대한 올바른 결론을 내리기 위해, 주어진 실험 저항성 데이터를 분석한다:
\begin{enumerate}[leftmargin=12pt, labelsep=4pt, itemsep=0pt, parsep=0pt, topsep=2pt]
\item \textbf{야생형 (control):} 100\% 저항성.
\item \textbf{g1 돌연변이체:} 75\% 저항성 (G1은 저항성에 약간 기여한다). \textit{[...]}
\end{enumerate}

\noindent\textbf{상호작용 분석:}
\begin{itemize}[leftmargin=12pt, labelsep=4pt, itemsep=0pt, parsep=0pt, topsep=2pt]
\item \textbf{Transcription Factor:} G2가 G1과 G3 모두의 upstream에 있는 transcription factor라면, G2를 knock-out할 경우 G1과 G3 모두의 발현이 막힌다. \textit{[...]}
\end{itemize}

\vspace{2pt}
\noindent\textbf{최종 답: (B)}
\end{mdframed}
\end{mdframed}
\end{minipage}
\caption{\textbf{GPQA-Diamond (Biology) 예시로, \sdrl 학습 전후 DiffusionGemma sampler에서의 모델 동작을 보여준다}. SFT 모델은 G2를 필수적인 transcription factor로 식별하고 G1/G3의 redundancy를 인식하는 \emph{올바른} 추론으로 시작하지만, 반복적인 token 루프 (숫자 `1')로 퇴화하면서 빈 최종 응답을 낸다. \sdrl 이후에는 모델이 분석을 깔끔하게 마무리하고 정답을 출력한다.}
\label{fig:gpqa-rl-example-biology}
\end{figure*}

\begin{figure*}[t]
\scriptsize
% === Full-width PROMPT ===
\begin{mdframed}[style=promptbox]
\noindent\textbf{프롬프트}
\vspace{2pt}\noindent
질량 $m$인 양자역학적 입자가 $(r, \theta)$의 함수로 주어지는 다음 퍼텐셜 안에서 2차원 운동을 한다:
%
\[
V(r, \theta) = \tfrac{1}{2}\,k\,r^2 + \tfrac{3}{2}\,k\,r^2\cos^2(\theta)
\]
%
에너지 스펙트럼을 구하라.

\vspace{2pt}
\noindent
\begin{tabular}{@{}ll@{}}
(A) & $E = (n_x + 3n_y + \tfrac{3}{2})\,\hbar\sqrt{k/m}$ \\[2pt]
(B) & $E = (2n_x + n_y + \tfrac{3}{2})\,\hbar\sqrt{k/m}$ \\[2pt]
(C) & $E = (3n_x + 2n_y + \tfrac{1}{2})\,\hbar\sqrt{k/m}$ \\[2pt]
(D) & $E = (2n_x + 3n_y + \tfrac{1}{2})\,\hbar\sqrt{k/m}$
\end{tabular}

\hfill\textbf{\textit{정답:} (B)}
\end{mdframed}
\vspace{2pt}
% === Two columns: Before \sdrlfootnotesize | After \sdrlfootnotesize ===
\noindent
\begin{minipage}[t]{0.49\textwidth}
% --- Before \sdrlfootnotesize ---
\begin{mdframed}[style=beforebox]
\noindent\textbf{\sdrlfootnotesize 이전} \hfill {\colorbox{beforebg}{\textcolor{beforeframe}{\textbf{$\boldsymbol\times$~오답 (비어 있음)}}}}

\noindent\textit{사고 과정 (일부 생략):}
\begin{mdframed}[style=thinkbox]
퍼텐셜은 극좌표 $(r, \theta)$에서 다음과 같이 주어진다:
$V(r, \theta) = \frac{1}{2}kr^2 + \frac{3}{2}kr^2\cos^2(\theta)$.
우리는 다음을 안다:

\quad $x = r\cos(\theta)$,

\quad $y = r\sin(\theta)$,

\quad $r^2 = x^2 + y^2$,

\quad $\cos^2(\theta) = x^2/r^2$.

\smallskip
이를 퍼텐셜 식에 대입하면:

\quad $V(x,y) = \frac{1}{2}k(x^2+y^2) + \frac{3}{2}kx^2 = 2kx^2 + \frac{1}{2}ky^2$

\textit{[...]}

$\omega_x$와 $\omega_y$ 값을 대입하면:

\textit{[...]}

계산된 에너지 스펙트럼은 $E = (2n_x + n_y + \tfrac{3}{2})\,\hbar\sqrt{k/m}$이며, 이는 선택지 (B)와 일치한다.

\medskip
\textit{[\ldots 답을 세 번 더 재검증하고, 모두 올바르게 (B)에 도달함 \ldots]}

\medskip
\noindent 선택지를 다시 읽어 보면
\quad .\qquad = ( = =xx\_nn\_\_33/ \qquad ))))
\smallskip
)))))))222\_\_\_\_y +//)) ))) ))))))))))))))2 2 2222\\
)))))))))))))))))))))))))))))))))))))))))))))))))\\
{\itshape [\ldots 최대 token 예산에 도달할 때까지 닫는 괄호가 반복됨.]}
\end{mdframed}
\vspace{2pt}
\noindent\textit{응답:}
\begin{mdframed}[style=thinkbox]
\textit{[비어 있음 --- 응답이 생성되지 않음]}
\end{mdframed}
\end{mdframed}
\end{minipage}%
\hfill
\begin{minipage}[t]{0.49\textwidth}
% --- After \sdrlfootnotesize ---
\begin{mdframed}[style=afterbox]
\noindent\textbf{\sdrlfootnotesize 이후} \hfill {\colorbox{afterbg}{\textcolor{afterframe}{\textbf{\checkmark~정답}}}}

\noindent\textit{사고 과정 (일부 생략):}
\begin{mdframed}[style=thinkbox]
퍼텐셜은 $V(r, \theta) = \frac{1}{2}kr^2 + \frac{3}{2}kr^2\cos^2(\theta)$이다.

\quad 좌표는 극좌표 $(r, \theta)$이다.

\quad $x = r\cos\theta$

\quad $y = r\sin\theta$

\quad $r^2 = x^2 + y^2$

\quad $\cos\theta = \frac{x}{r}$

\quad $V(r, \theta) = \frac{1}{2}k(x^2+y^2) + \frac{3}{2}k(x^2+y^2)(\frac{x}{r})^2$

\textit{[\ldots]}

Hamiltonian은 두 개의 1차원 조화진동자로 분리될 수 있다:

\quad $H = H_x + H_y$

\textit{[\ldots]}

1D 조화진동자 $H = -\frac{\hbar^2}{2m}\frac{d^2}{dx^2} + \frac{1}{2}m\omega^2x^2$의 에너지 고유값은 다음과 같다:
    $E = (n+\frac{1}{2})\hbar\omega$, 여기서 $n = 0, 1, 2, ...$이다.

\textit{[\ldots]}

계산된 $E = (2n_x+n_y+3/2)\hbar\sqrt{k/m}$는 선택지 (B)와 일치한다.
   
\textit{[\ldots 계산을 두 번 검증하고, 선택지 (B)와의 일치를 확인함.]}
\end{mdframed}
\smallskip
\noindent\textit{응답 (발췌):}
\begin{mdframed}[style=thinkbox]
주어진 퍼텐셜 $V(r, \theta) = \frac{1}{2}kr^2 + \frac{3}{2}kr^2\cos^2(\theta)$에서 질량 $m$인 양자역학적 입자의 에너지 스펙트럼을 구하기 위해, 먼저 퍼텐셜을 Cartesian 좌표 $(x, y)$로 변환한다.

1.  \textbf{퍼텐셜을 Cartesian 좌표로 변환한다:}
    극좌표와 Cartesian 좌표의 관계는 다음과 같다:

\textit{[\ldots]}

    % Substituting these into the potential formula:
    % $V(r, \theta) = \frac{1}{2}k(x^2+y^2) + \frac{3}{2}k(x^2+y^2)\left(\frac{x}{r}\right)^2$

% \textit{[\ldots]}

2.  \textbf{Hamiltonian을 식별한다:}
    2차원에서 입자의 Hamiltonian은 다음과 같다: \textit{[\ldots]}

    % This is the sum of two independent one-dimensional harmonic oscillators, $H = H_x + H_y$.

3.  \textbf{각 진동자의 에너지 고유값을 계산한다:}
    1D 조화진동자 $\frac{1}{2}m\omega^2x^2$의 에너지 준위는 $E = (n+\frac{1}{2})\hbar\omega$이다.

\textit{[\ldots]}

4.  \textbf{전체 에너지를 구한다:}

\textit{[\ldots]}

\quad $E = (2n_x + n_y + \frac{3}{2})\hbar\sqrt{\frac{k}{m}}$

\textbf{최종 답: (B)}
\end{mdframed}
\end{mdframed}
\end{minipage}
\caption{\textbf{GPQA-Diamond (2D harmonic oscillator energy spectrum) 예시로, \sdrl 학습 전후 DiffusionGemma sampler에서의 모델 동작을 보여준다}. SFT 모델은 Cartesian 좌표로의 변환, 분리 가능한 Hamiltonian의 식별, 올바른 에너지 스펙트럼 계산이라는 \emph{정확한} 유도로 시작하지만, 반복적인 token 루프 (닫는 괄호)로 퇴화하면서 빈 최종 응답을 낸다. \sdrl 이후에는 모델이 같은 유도를 깔끔하게 마무리하고 정답을 출력한다.}
\label{fig:gpqa-rl-example-quantum}
\end{figure*}
